\section{Discussion}\label{sec:discussion}

\subsection{Principal findings}

\pending{Summarise the paired differences for H1--H4 by model, with their confidence intervals, without
ranking models and without extrapolating beyond this cohort and schema.} At the current data cut, the
measured results concern the reference rather than the models. Pairwise agreement between reviewers on
the 174 boolean variables was substantial overall ($\kappa=0.76$), but agreement on the presence of a
finding (positive specific agreement 0.79) was lower than on its absence, and agreement was lowest for
complications and infections, the blocks whose manual rules require the most interpretation (the
infection cascade, colonisation versus infection, the severity threshold of delirium). Any model
evaluated against this reference inherits that uncertainty, which bounds the agreement that can be
expected and argues for reporting results by block.

\subsection{Why decomposition may help}

Four mechanisms favour the harness, and each maps onto a measured outcome. First, each call produces a
short output, so the risk of reaching the token limit falls (H1). Second, validation is local: an error
invalidates one group of at most 16 variables rather than the whole form, and the retry receives a
short, specific error list (H1, retry recovery). Third, each prompt carries only the definitions of the
requested variables, so the model does not need to keep 199 definitions in view while generating
(H2, H3). Fourth, the harness is operationally robust: signed checkpoints allow a job interrupted by the
scheduler to resume without repeating validated groups, which matters on a shared cluster with wall-time
limits (H4).

\subsection{Why decomposition may hurt}

Decomposition also has costs and failure modes that this design is built to expose rather than to hide.
The note is repeated in all 18 prompts, so input tokens and prefill time grow roughly with the number of
groups (H4). Related variables are answered in separate calls: the infection cascade spans three groups, and
51 of the 137 parent--child relations of the form (affecting 44 variables) and 3 of its 7 mutual-exclusion
pairs cross group boundaries. The per-group retry cannot
repair such inconsistencies, which are detected only when the groups are merged; the harness therefore
may produce more case-level relation errors than a single call that sees the whole form. Each call is
also asked, in isolation, to look for a small set of findings, which could bias it toward positive
answers for rare variables, as observed with category-specific prompts in earlier
work~\cite{agrawal2022fewshot}; the error taxonomy includes an over-calling category to detect this. Finally,
the partition is mechanical, consecutive variables within a block, and was not optimised; a different
grouping could change the results in either direction.

\subsection{What a shared case model can and cannot add}

Arm~C is designed to keep the advantages of decomposition (short outputs, local validation, batching) while
restoring what independent calls lose: a common reading of the episode. Its expected benefit is concentrated
where the pilot and the literature locate the difficulty. In a two-stage pipeline on psychiatric discharge
summaries, an intermediate extraction of events and times improved over direct extraction mainly for the most
inferential feature, the onset of the first episode~\cite{psych2stage2026}; in cancer registry abstraction, a
multi-agent pipeline improved context-dependent variables such as tumour grade but not highly structured ones
such as TNM stage~\cite{cancerregistrymas2026}. H6 therefore expects gains in infections, organ failure,
complications and ICU discharge rather than in medical history, and larger gains in poorly documented notes,
where information is scattered across sections or must be inferred. The case model also makes explicit whether
an ICU stay occurred at all, which in Arm~B each group decided separately. The same studies caution that
decomposition redistributes performance rather than improving it uniformly: in the cancer registry study
short, standardised fields such as metastasis stage lost recall~\cite{cancerregistrymas2026}, and in the
psychiatric study the intermediate layer did not help simple counts~\cite{psych2stage2026}. Its main risk is
propagation. In document-level extraction, context injected into every chunk outperformed both no context and
a single static summary, but an erroneous shared structure performed worse than no structure at all, with F1
falling from 0.80 to 0.45~\cite{hisccg2026}; the designers of the first blackboard system similarly reported that
focusing the search on low-level hypotheses whose credibility ratings were insufficiently accurate was
ineffective and harmful~\cite{erman1980hearsay}. This is why the case model must cite the note, the groups must give priority to
the note, the case model is audited by a clinician, and propagated errors are a category of the error
taxonomy. \pending{interpret H5 and H6.}

\subsection{Evidence by line identifier}

Citing evidence by line identifier and copying the text in software removes one failure mode, the
fabricated or paraphrased quotation, in both arms. It does not guarantee that the cited line supports the
answer, and a line of a PDF-derived text can be too coarse (several findings on one line) or too fine (a
sentence split across visual lines, or a table flattened into lines). For this reason evidence fidelity
is measured as identifier validity, coverage and clinician-judged support, not as literal matching, which
holds by construction. \pending{interpret H2 results.}

\subsection{Relation to previous work}

The direct comparisons available in the literature point in different directions and should not be read
as a general rule. Splitting a note into sentences raised F1 substantially over a single whole-document call
for dense entity extraction, at about twice the GPU time~\cite{llmie2025}; a prompt graph of sequential
subtasks gave the largest mean gain over zero-shot prompting in a multilingual benchmark, at roughly eight
times the time per report of few-shot prompting~\cite{spaanderman2025structured}; and sequential prompting did
not outperform few-shot prompting for timeline extraction with a 7B model~\cite{fornasiere2024medical}. None of
these studies decomposed a fixed schema of this size by output groups, so they motivate the hypotheses but do
not predict the result. \pending{relate the measured differences to these findings.}

Building an intermediate representation before the final extraction also has precedents. Clinical pipelines
have extracted events and temporal information before predicting chart-level features~\cite{psych2stage2026}
or reconstructing patient timelines~\cite{clinicirca2026}, and a pathology workflow passed the metadata of a planning step to parallel field extractors, although
without testing whether that shared context helped~\cite{promptplanextract2026}. Outside medicine, blackboard architectures
coordinated independent specialists through a shared data structure~\cite{erman1980hearsay,nii1986blackboard},
an idea recently revived for multi-agent LLM systems~\cite{blackboardllm2025,blackboardmas2025}, and
Skeleton-of-Thought generates a plan first and expands its parts in parallel~\cite{skeletonofthought2024}, the
same computational pattern as Arm~C. What these works do not provide is a paired comparison, on the same
patients and schema, between a single call, independent decomposed calls and decomposed calls that share a
verified representation of the case.

Two observations from earlier work are specific enough to be tested here. Agrawal et al.\ found that prompts
specialised for one category tended to produce an answer even when none was present~\cite{agrawal2022fewshot};
if the harness behaves similarly, it should show more positive answers for rare variables than the monolithic
call, which the over-calling category of the error taxonomy and the per-block positive rates are designed to
detect. Hein et al.\ observed more major errors in reports with many items to extract and argued that success
depended on the clarity of the instructions~\cite{iterativerefine2025}; the instruction text here was derived
from the annotators' manual rather than refined against the reference, which avoids overfitting the prompts
to the reference but may leave ambiguities that only an error analysis will reveal.

Other design choices in the literature address the same problems by different means. Grammar-constrained
decoding guarantees well-formed JSON~\cite{privacy2024structured} and would remove parse failures from all
arms, which is why parse success is reported separately from content validity. Entity-anchored retrieval
reduces the input sent to each call~\cite{clinicalentityretrieval2024}; in the harness the full note is repeated
in every call, so input-side decomposition is a natural complement to the output-side decomposition studied
here. Line-numbered input for evidence has been used before~\cite{fornasiere2024medical}, and criterion-level
decisions with sentence locations have been used for trial matching~\cite{trialgpt2024}; the difference here is
that identifiers are validated, the quotation is copied by software and the support of the cited lines is
assessed by clinicians.

\subsection{Limitations}

The study has a single centre, a single language and a single form, and 50 summaries; results may not
transfer to other hospitals, note styles or schemas. The cohort was not screened for the presence of an ICU
stay, and about a quarter of the summaries may not document one; the agreement analyses are restricted to
eligible summaries, which reduces the sample. Arm~C and hypotheses H5 and H6 were added after the pilot; they
were specified before any comparison with the human reference, but they were motivated by pilot observations.
The evaluation engine may differ from the pilot engine in precision (FP8 instead of NF4) and in batching, so the
pilot is not a calibration of the evaluation run. The human reference was incomplete at the data cut,
reviewers did not use the ? option, and no adjudicated consensus exists; individual references attenuate
absolute agreement, although they affect both arms equally. Positive findings are infrequent and 40
boolean variables were never marked positive, so per-variable estimates are unstable or undefined. The
model input is PDF-derived text grouped into visual lines, which can disorder tables or columns;
de-identification relied on pattern-based redaction, which masked parts of clinical words in 3 of 655
redacted segments, including a sedative drug name. The two arms necessarily differ in their output budget
and in the length of their prompts; the pre-specified budgets and the sensitivity analysis at an equal
per-call budget address this but cannot remove it. Greedy decoding on GPUs may not be bit-reproducible,
especially when prompts are batched and kernels are not batch-invariant; the batch-invariant mode, the
equivalence test and the repeatability subset address this, but only on the cases tested. Three models were evaluated, each with one run; conclusions about other models or sampling
settings would require new runs. Finally, the paired design isolates the calling strategy but not its
components (grouping, retries, per-group validation); the ablations without retries address only one of
them.

\subsection{Implications}

\pending{State what the measured trade-off implies for the design of extraction pipelines for long clinical
forms, limited to the measured outcomes.} Independently of the results, the protocol illustrates reporting
practices that make such comparisons auditable: frozen inputs with hashes, evidence that cannot be
fabricated, failures that are never converted into negative findings, and a reference whose agreement and
missingness are reported before it is used.
